\label{sec:fairness}

Fairness in apportionment is multidimensional.
Three criteria have dominated the academic and legal literature:
the \emph{population pair test} (which method is fairest to each pair
of states?), the \emph{quota rule} (is each state within 1 of its
exact proportional entitlement?), and \emph{paradox immunity}
(does the method produce counterintuitive results when House size
or populations change?).
No method satisfies all three simultaneously --- the Balinski-Young
impossibility theorem guarantees this \citep{balinski1982fair}.

\subsection{The Population Pair Test}

For any pair of states $(i,j)$, the per-capita representation ratio
is $R_{ij} = (p_i/s_i) / (p_j/s_j)$.
In a perfectly proportional system, $R_{ij} = 1$ for all pairs.
The \emph{population pair test} asks: does the method minimise
the maximum deviation of $R_{ij}$ from 1 over all pairs?

Huntington-Hill minimises the maximum \emph{relative} deviation
$\max_{i \neq j} |R_{ij} - 1|/\min(R_{ij}, 1/R_{ij})$, which Huntington
\citeyearpar{huntington1928method} characterised as minimising the
relative difference in per-capita representation.
Webster minimises the maximum \emph{absolute} deviation
$\max_{i \neq j} |p_i/s_i - p_j/s_j|$.
The two criteria agree for most pairs but can disagree for pairs of
states with very different populations (e.g., California vs.\ Wyoming),
where the choice between relative and absolute deviation produces
different seat assignments.

\subsection{The Quota Rule}

The quota rule requires that each state receive either $\lfloor q_i \rfloor$
or $\lceil q_i \rceil$ seats, where $q_i = p_i \times H / N$ is its
exact proportional quota.

\begin{center}
\begin{tabular}{ll}
\toprule
Method & Quota rule satisfied? \\
\midrule
Hamilton    & Yes (by construction) \\
Huntington-Hill & Yes for 2020 Census; not guaranteed in general \\
Webster     & Yes for 2020 Census; not guaranteed in general \\
Adams       & Always at upper quota ($s_i \geq \lceil q_i \rceil$) \\
Jefferson   & Always at lower quota ($s_i \leq \lceil q_i \rceil$); may violate upper \\
\bottomrule
\end{tabular}
\end{center}

Hamilton's quota compliance is exact by construction.
Huntington-Hill and Webster satisfy the quota rule in practice for
all known U.S.\ census data but can in principle violate it for
extreme population distributions.
The Balinski-Young theorem \citep{balinski1982fair} shows that
satisfying the quota rule for all possible population vectors
requires a quota method, which is incompatible with paradox immunity.

\subsection{Paradox Immunity}

The three classical apportionment paradoxes are:
\begin{enumerate}
  \item \textbf{Alabama Paradox}: Increasing $H$ by 1 causes some state
    to lose a seat.
  \item \textbf{Population Paradox}: State $A$'s population grows faster
    than state $B$'s between censuses, yet $A$ loses a seat to $B$.
  \item \textbf{New-States Paradox}: Adding a new state with proportional
    seats causes existing states to exchange seats.
\end{enumerate}

Divisor methods are immune to all three paradoxes (proved in J.5).
The key insight is that in a priority-queue divisor method,
awarding additional seats can only go to states with the highest
priority --- a state that currently holds seat $s$ holds it because
its priority $P_i(s-1)$ exceeded all competitors' priorities at that
step.
Increasing $H$ adds more seats through the priority queue;
no existing seat allocation is reconsidered.
Hamilton is susceptible to all three paradoxes because its
fractional-remainder ranking is non-monotone in $H$ and population.

\begin{center}
\begin{tabular}{lccc}
\toprule
Method & Alabama & Population & New-States \\
\midrule
Huntington-Hill & Immune & Immune & Immune \\
Webster         & Immune & Immune & Immune \\
Adams           & Immune & Immune & Immune \\
Jefferson       & Immune & Immune & Immune \\
Hamilton        & Susceptible & Susceptible & Susceptible \\
\bottomrule
\end{tabular}
\end{center}

\subsection{Bias Direction}

The four divisor methods lie on a spectrum from most small-state-favoring
(Adams) to most large-state-favoring (Jefferson), with
Huntington-Hill and Webster in between:

\begin{center}
Adams $\succ$ Huntington-Hill $\succ$ Webster $\succ$ Jefferson
\end{center}

\noindent
(Where $A \succ B$ means $A$ gives systematically more seats to small
states relative to $B$.)
This ordering follows from the rounding thresholds:
Adams always rounds up (giving small states the benefit of fractional
seats), Jefferson always rounds down (denying fractional seats),
Huntington-Hill rounds at the geometric mean (slightly above 0.5,
favouring small states), and Webster rounds at 0.5 (neutral between
large and small states).

For the 2020 Census, small-state bias manifests most visibly in the
treatment of states near the 1-to-2 seat boundary.
Wyoming, Vermont, and Alaska all have quotas between 0.75 and 1.0:
they receive 1 seat under all five methods because the constitutional
floor overrides rounding, but their \emph{effective} over-representation
relative to strict proportionality is highest under Adams and lowest
under Jefferson.
